% Independent proof review accompanying the machine-readable audit receipt.
% The principal construction is in division_norm_compiler_20261106_v1.tex.
\section{Independent verification of the cyclic-norm specifications}
\label{sec:independent-norm-audit-phase3}

The independent checker
\texttt{verify\_division\_norm\_specs\_independent\_20261106\_v1.py}
imports neither the producer nor an earlier checker.  It compares the
complete residual lists with pinned source bytes, constructs every
cyclic-matrix entry by iterating the shift operator, verifies its formal
Galois conjugacy, and recomputes the degree and coefficient bounds.
Its receipt is
\texttt{division-norm-independent-verification\_20261106\_v1.json}.
The statements below explain the algebra certified by those finite checks.
The determinants themselves are not expanded, and this audit does not
replace the arithmetic proof identifying a source's projection.

\subsection{The determinant and its exact degree}

Let $d,q,p$ be one of $(36,37,2)$ and $(37,149,5)$.  Put
$L=\mathbb Q(\zeta_q)^H$, where $H$ is the subgroup of order
$(q-1)/d$ of $(\mathbb Z/q\mathbb Z)^*$, and choose the generator
$\sigma$ determined by the recorded primitive root.  The periods
$\eta_i=\sum_{h\in H}\zeta_q^{g^ih}$ partition the nonzero
exponents modulo $q$.  Their linear independence follows because
a rational relation gives a polynomial of degree at most $q-1$
with zero constant coefficient divisible by
$1+T+\cdots+T^{q-1}$, hence the zero polynomial.  Their sum is
$-1$, so $1,\eta_1,\ldots,\eta_{d-1}$ is also a basis.

In the cyclic algebra $A=(L/\mathbb Q,\sigma,p)$, write
$a_j=Z_{j,0}+\sum_{i=1}^{d-1}\eta_iZ_{j,i}$ and
$a=\sum_{j=0}^{d-1}a_ju^j$.  Let
\[
 D(a)=\operatorname{diag}(a,\sigma(a),\ldots,\sigma^{d-1}(a)),
 \qquad Ue_j=e_{j-1}\ (j>0),\quad Ue_0=pe_{d-1}.
\]
Then $U^d=pI$ and $UD(a)=D(\sigma(a))U$.  Thus the split
matrix of $a$ is $M=\sum_jD(a_j)U^j$.
The checker verifies every entry of this expression and the identity
\[
                            \sigma(M)=UMU^{-1}
\]
as a formal linear form in all the $Z_{j,i}$.  It follows that
$\det M$ is fixed by a generator of the Galois group and has
rational coefficients.  Its coefficients are algebraic integers,
so they are integers.  Reducing modulo the cyclotomic polynomial
therefore leaves only its constant coefficient in $\zeta_q$.
The specified coefficient-extraction operation is the actual reduced
norm, with no cyclotomic variable existentially quantified.

In each of the four supplied residual tables the input $X_0$ occurs
only in the term $+X_0^2$ of residual $F_0$.  The first coefficient
basis element is one, so the substituted matrix is
\[
                           X_0^2 I+M_0,
\]
where $M_0$ is independent of $X_0$.  Consequently its determinant
has coefficient one at $X_0^{2d}$.  Every residual has total degree
at most two, so the determinant has total degree exactly $2d$.
This also proves that the resulting integer polynomial is primitive.

\subsection{Division at the specified primes}

At $p$, the extension $L_p/\mathbb Q_p$ is unramified of degree
$d$: the checker computes the order of $pH$ in the quotient group
and finds $d$.  With $v(p)=1$, the function
\[
 w\left(\sum_{j=0}^{d-1}a_ju^j\right)
       =\min_{a_j\ne0}\left(v(a_j)+\frac jd\right)
\]
has a unique smallest term.  The product of the two unique smallest
terms is strictly smaller than every other product term; replacing
$u^d$ by $p$ preserves this weight.  Thus $w(ab)=w(a)+w(b)$,
so the local algebra has no zero divisors.  A finite-dimensional
algebra over a field with no zero divisors is a division algebra.
This proves division for both recorded rational cyclic algebras.

For use over $\mathbb Q(i)$, the degree-36 construction also needs
a division place at 37: its local index at two alone does not
establish degree-36 division after a quadratic extension.

\begin{lemma}\label{lem:norm36-local37-independent}
The algebra
\[
 \bigl(\mathbb Q_{37}(\zeta_{37})/\mathbb Q_{37},
          \ \sigma(\zeta_{37})=\zeta_{37}^2,\ 2\bigr)
\]
is a division algebra of degree 36.  Hence the specified rational
degree-36 algebra remains division over $\mathbb Q(i)$.
\end{lemma}
\begin{proof}
Put $\pi=\zeta_{37}-1$.  The polynomial
$\Phi_{37}(1+T)=((1+T)^{37}-1)/T$ is Eisenstein at 37.
Thus the extension is totally ramified of degree 36, with
uniformizer $\pi$ and residue field $\mathbb F_{37}$, and
\[
                  \frac{\sigma^j(\pi)}{\pi}\equiv2^j\pmod\pi.
\]
The checker verifies that two has multiplicative order 36 modulo
37.  The polynomial $T^{36}-2$ is irreducible over
$\mathbb F_{37}$: any root $\alpha$ satisfies
$\alpha^{37}=2\alpha$, and hence
$\alpha^{37^m}=2^m\alpha$.  Its Frobenius orbit therefore has
length exactly 36, which is the degree of the polynomial.
Thus
\[
                 R=\mathbb F_{37}[t]/(t^{36}-2)
\]
is a field.

Normalize $v_\pi(\pi)=1$.  For a nonzero element of the cyclic
algebra define
\[
               w\left(\sum_{j=0}^{35}a_ju^j\right)
                       =\min_jv_\pi(a_j).
\]
Let $x=\pi^m\sum b_ju^j$ and $y=\pi^n\sum c_ju^j$, with
all $b_j,c_j$ integral and at least one unit in each list.
Their leading residues $B(t),C(t)$ are nonzero elements of $R$.
Commuting $\pi^n$ past $u^j$ shows that the leading residue of
$\pi^{-m-n}xy$ is
\[
                              B(2^nt)C(t).
\]
The substitution $t\mapsto2^nt$ is an automorphism of $R$,
since $2^{36n}=1$ in the residue field; this remains true for
negative $n$.  The displayed product is therefore nonzero.
Thus $w(xy)=m+n$, and the algebra has no zero divisors.  It is
division by finite dimensionality.

Finally $6^2\equiv-1\pmod{37}$, and the derivative of
$T^2+1$ at six is nonzero modulo 37.  Hensel's lemma shows
that 37 splits in $\mathbb Q(i)$ and its completions there
are $\mathbb Q_{37}$.  The base-changed algebra is division
at such a completion.  A zero divisor over $\mathbb Q(i)$
would remain a zero divisor after that scalar extension, so
none exists.
\end{proof}

\begin{lemma}\label{lem:coprime-division-persistence-independent}
Let $D$ be a central division algebra of degree $d$ over a field
$k$, and let $K/k$ be a finite extension of degree $n$.
If $\gcd(d,n)=1$, then $D\otimes_kK$ is division.
\end{lemma}
\begin{proof}
Write $D\otimes_kK\simeq M_s(D')$, where $D'$ is central
division over $K$.  Degree comparison gives $\deg D'=d/s$.
A simple left module for this matrix algebra has $K$-dimension
$s(d/s)^2=d^2/s$, and hence $k$-dimension $nd^2/s$.
Restrict its action to $D$.  A finite-dimensional module over
the division algebra $D$ is a vector space over $D$, so its
$k$-dimension is a multiple of $\dim_kD=d^2$.
Consequently $s\mid n$.  Also $s\mid d$, so coprimality
forces $s=1$.
\end{proof}

This module argument is required for a general composite degree.
Coprimality of $d$ with the global extension degree does not imply
that one particular completion has degree coprime to $d$.
For degree 37 the simpler place-by-place argument in the principal
construction is valid, because 37 is prime.

\subsection{Coefficient bounds and finite-algebra degree bounds}

Each period has complex absolute value at most $(q-1)/d$.
Every entry of $M$ therefore has coefficient norm at most
$(q-1)p$, so
\[
                   \|\det M\|_1\le d!\bigl((q-1)p\bigr)^d.
\]
The exact integer checks give
\[
                     36!72^{36}<2^{361},\qquad
                     37!740^{37}<2^{496}.
\]
The maximum coefficient norm of any recorded residual, with the
fixed integer $\kappa$ treated as a formal coefficient symbol,
is at most $2^{511}$.  In the two symbolic-constant sources its
exponent is at most one.  Since $3^5<2^8$,
\[
             \log_2\kappa=2^{28}\log_2 3
                         <\frac85\,2^{28}<429496730.
\]
Substitution and submultiplicativity of coefficient norm therefore
give the independently verified bounds
\[
\begin{array}{c|c|c}
 \text{total coordinates}&\text{exact degree}&H(P)\ \text{at most}\\\hline
 1315&72&2^{15461901037}\\
 1339&72&2^{18757}\\
 1367&74&2^{15891398413}\\
 1391&74&2^{19403}
\end{array}
\]

For the finite-algebra outputs, the checker independently follows
the coefficient DAG to verify generator weights, triangular
dependencies, and companion-operator norm inequalities.  It
recounts all parameter, quadratic-mask, separating-parameter,
and degree cases.  If $p$ denotes the number of quartic point
generators, $q$ the number of quadratic generators, $m=p+q$,
and $B$ the number of finite parameter branches, this count is
\[
 B\sum_{k=0}^q\binom qk\,4^p2^k
        \left(1+\frac{m(4^p2^k)(4^p2^k+1)}2\right).
\]
With rank bound $n$ and $c$ closed residuals, each branch has at
most $n+c$ equations.  A norm of degree
$d=\lceil\sqrt{n+c}\rceil$ replaces the sum of squares, and
multiplies the existing branch-degree bound by $d/2$.
The checker uses the exact factor count and unrounded branch
majorant before taking a binary logarithm.  This matters: first
rounding the old total bound and the multiplier separately can
lose a bit.  The resulting verified estimates are
\[
\begin{array}{c|c|r|c}
 \text{source}&\text{coordinates}&d&\deg P\ \text{at most}\\\hline
 \text{new finite anchor}&41&16385&2^{9566686}\\
 \text{new semilocal}&43&23171&2^{250}\\
 \text{previous finite anchor}&43&23171&2^{9566693}\\
 \text{previous semilocal}&45&32769&2^{257}
\end{array}
\]
The receipt retains the exact integer multipliers and the exact
slack in the last rounding step.  These statements concern
encoding a specified finite source over a chosen number field.
They assert neither a full coefficient expansion nor that the
base-changed source defines that field's ring of integers.
